% Clase 2 — Planeación: Fundamentos
% Lectura: AIMA Secciones 11.1–11.2
\section{Planeación: Definición y Algoritmos Generales}

\subsection{¿Qué es la Planeación?}

Un agente planificador construye una \textbf{secuencia de acciones} para lograr un objetivo en un entorno modelado.

\begin{definition}[Problema de Planeación STRIPS/PDDL]
\begin{itemize}
    \item \textbf{Estado}: conjunción de fluentes (literales positivos/negativos).
    \item \textbf{Acción}: definida por precondiciones y efectos.
    \item \textbf{Meta}: conjunción de fluentes que deben ser verdaderos.
\end{itemize}
\end{definition}

\subsection{Representación PDDL}

\begin{lstlisting}[language=, caption={Ejemplo de acción en PDDL}]
(:action Fly
 :parameters (?p - Plane ?from ?to - Airport)
 :precondition (and (At ?p ?from) (Plane ?p)
                    (Airport ?from) (Airport ?to))
 :effect (and (At ?p ?to) (not (At ?p ?from))))
\end{lstlisting}

\subsection{Algoritmos de Búsqueda para Planeación}

\begin{itemize}
    \item \textbf{Búsqueda hacia adelante} (\textit{progression}): empieza en el estado inicial y aplica acciones.
    \item \textbf{Búsqueda hacia atrás} (\textit{regression}): empieza desde la meta y aplica acciones inversas.
    \item \textbf{SATPLAN}: codifica el problema de planeación como SAT y usa un solucionador SAT.
\end{itemize}
